\scanpage{27}
% archive sheet number: 26

\deleted{
\diagram{\begin{tikzcd}[ampersand replacement=\&]
 A\times B\arrow[r,"s_{A,B}"]\arrow[d,"f\times g"] \& B\times A\arrow[d,"g\times f"]\\
 A'\times B'\arrow[r,"s_{A',B'}"] \& B'\times A'
\end{tikzcd}}
\[
\begin{aligned}
(g\times f)_*(s(x\otimes y))&=(g\times f)_*(y\otimes x)(-1)^{xy}
=\unsure{\{}g_*(y)\otimes f_*(\unsure{y|x})(-1)^{xy+\text{\textcircled{$fx$}}}\\
s(f\times g)_*(x\otimes y)&=s(f_*(x)\otimes g_*(y))(-1)^{\text{\textcircled{$gx$}}}
=g_*(y)\otimes f_*(\unsure{x})(-1)^{\text{\textcircled{$gx$}}+(x+f)(y+g)}
\end{aligned}
\]
\[
\begin{aligned}
&(-1)^{\deleted{xy}+fx+\deleted{gx}+\deleted{gx}+fy+\deleted{xy}+fg}
\qquad\deleted{=(-1)^{(f+x)(g+y)}}\\
&\unsure{\shortparallel}\qquad\\[-.4em]
&(-1)^{f(x+y+g)}\qquad(-1)^{\deleted{\unsure{xy+}}fy+gx+fg}
\end{aligned}
\]
}
\noindent\rule{.91\linewidth}{.4pt}

\uline{A}

1) Catégorie \inserted{$V^\circ$} des alg. de Poincaré $A$ sur $k$ : alg. gr. \unsure{comm.} proj. de t. f. à degrés $\geq0$, munie d'un hom. lin. $k_A\colon A\mathbin{-}k$, de degré $d^\circ(A)$ dépendant de $A$ tel que $(x,y)\longmapsto k_A(xy)$ soit une forme \inserted{strictement} non dégénérée. \deleted{\unsure{pour} \unsure{strict}}

Les hom. \inserted{de $V^\circ$} : ce sont les hom. d'algèbres gr. \unsure{degrés|unités}, sans plus (donc il n'est pas question dans la définition de $V$ hom des $k_A$)

\margin{$\unsure{Z|2}$}
\uline{NB} Un hom de cette catégorie n'est pas néc. compatible avec \inserted{les} $k_A,k_B$

2) Foncteur contravariant $\varphi\colon (V^\circ)^\circ=V\longrightarrow\deleted{\unsure{M}}\unsure{\mathrm{Mod\,gr.d}|\mathrm{Modgrad}}(k)$

a) \unsure{Sur|sur} les objets : $\varphi(A)=A$ où on oublie la multiplication

b) Sur les morphismes : soit $f^*\colon B\mathbin{-}A$, on pose $\varphi(f^*)=f_*\colon A\longrightarrow B$ défini par la condition que
\[
k_A(f^*(b)a)=k_B(bf_*(a))
\]
On vérifie $\varphi(g^*f^*)=\varphi(\unsure{f^*|f})\varphi(g^*)$, $\varphi(\mathrm{id})=\mathrm{id}$, et formule de proj.
\[
f_*(a)b=f_*(af^*b)
\]
enfin \deleted{$k_B(f_*(a))$} $k_B f_*=k_A$
